% BigAlpha 2026 — AB Test Report (Revised after Peer Review)
% ACM Conference Format (acmart)
\RequirePackage{amsmath,amssymb}
\let\Bbbk\relax
\documentclass[sigconf,10pt]{acmart}

\usepackage{booktabs}
\usepackage{multirow}
\usepackage{graphicx}
\usepackage{longtable}
\usepackage{array}

% ACM Metadata
\acmConference[ACML 2026]{Asian Conference on Machine Learning}{Dec 2026}{Beijing, China}
\acmYear{2026}
\setcopyright{none}
\settopmatter{printacmref=false}

\title{BigAlpha 2026: A Systematic Benchmark of 20+ Deep Learning Models for End-to-End Stock Return Prediction}

\author{Zhenhao Li}
\affiliation{%
  \institution{XJTLU / University of Liverpool}
  \city{Suzhou}
  \country{China}
}
\email{omitted@blind.review}

\begin{document}

\begin{abstract}
We present a comprehensive benchmark infrastructure of 23 deep learning and gradient-boosted models for end-to-end stock return prediction, built for the BigAlpha 2026 competition. Our model zoo spans five architectural families: (1) financial foundation models (Kronos~\cite{kronos2025}), (2) time-series transformers (Autoformer~\cite{wu2021autoformer}, Informer~\cite{zhou2021informer}, TimesNet~\cite{wu2023timesnet}, PatchTST~\cite{nie2023patchtst}), (3) recurrent architectures (LSTM~\cite{hochreiter1997long}, GRU~\cite{cho2014learning}), (4) Bayesian-optimized tree models (LightGBM~\cite{ke2017lightgbm}, XGBoost~\cite{chen2016xgboost}), and (5) multi-agent reinforcement learning (CTDE-MARL~\cite{ctde2026acml}). All models operate on raw OHLCVA features with zero feature engineering. We implement TPE-based Bayesian hyperparameter optimization~\cite{bergstra2011algorithms,snoek2012practical} via Optuna~\cite{akiba2019optuna} across 10 model families. A local baseline test on 6 models (10 A+H stocks, CPU, 20 epochs) confirms that daily data alone is insufficient for cross-sectional prediction (max IC = +0.038). We provide a three-phase competition strategy targeting IC $>$ 0.08 through IC-weighted ensemble of Bayesian-optimized models on GPU with minute-frequency data.
\end{abstract}

\maketitle

%=============================================================================
\section{Introduction}
%=============================================================================

The BigAlpha 2026 Global University League, organized by BigQuant, challenges participants to build end-to-end models predicting cross-sectional stock returns from raw market data. The End-to-End Price Prediction track imposes: (1) zero feature engineering, (2) input fields $\le 100$, (3) parameters $\in [10^5, 10^8]$, (4) from-scratch training, and (5) $\le 3$ hours GPU inference. Scoring weights: Rank IC (25\%), ICIR (25\%), Sharpe (25\%), and Stress robustness (25\%).

We present the model zoo and training infrastructure built for this competition. Drawing on our prior work on CTDE-MARL for portfolio management~\cite{ctde2026acml} and recent time-series architectures, we implement 23 models and apply TPE-based Bayesian optimization to each. This report documents: (1) the model design (§2), (2) Bayesian optimization search spaces (§3), (3) local baseline experiments (§4), and (4) competition submission strategy (§5).

\paragraph{Contributions.} (1) A unified 23-model benchmark codebase for E2E stock prediction. (2) TPE-based Bayesian optimization covering 10 model families with 7--10 parameter search spaces. (3) Local baseline experiments validating the necessity of GPU and minute-frequency data. (4) A three-phase competition strategy informed by model characteristics.

\paragraph{Scope and Limitations.} This report documents the \emph{infrastructure} and \emph{local baseline} — BigQuant GPU experiments with Bayesian optimization are in progress and will be reported separately. Results are currently limited to CPU-based daily-frequency evaluation; conclusions about relative model performance require GPU experiments. Missing baselines include iTransformer, DLinear, Mamba, and TimesFM which will be added in future work.

%=============================================================================
\section{Model Zoo}
%=============================================================================

All models accept raw 6-dimensional OHLCVA sequences (open, high, low, close, volume, amount) with configurable lookback windows. Each model family has a corresponding Bayesian optimizer (Table~\ref{tab:models}).

\subsection{Financial Foundation Models}

\textbf{Kronos}~\cite{kronos2025} is a decoder-only Transformer pre-trained on 12B+ K-line records from 45+ exchanges. The BSQ tokenizer converts OHLCVA into discrete tokens, naturally satisfying the zero-feature-engineering constraint. Variants: Kronos-mini (4.1M), Kronos-small (24.7M), Kronos-base (102.3M).

\subsection{Time-Series Transformers}

\textbf{Autoformer}~\cite{wu2021autoformer} uses auto-correlation and progressive decomposition. \textbf{Informer}~\cite{zhou2021informer} uses ProbSparse attention ($O(L \log L)$). \textbf{TimesNet}~\cite{wu2023timesnet} transforms 1D$\to$2D via FFT then applies Inception convolutions. \textbf{PatchTST}~\cite{nie2023patchtst} segments sequences into patches with channel-independent Transformers. Standard \textbf{Transformer}~\cite{vaswani2017attention} serves as baseline.

\subsection{Recurrent Architectures}

\textbf{LSTM}~\cite{hochreiter1997long} and \textbf{GRU}~\cite{cho2014learning} encode 60-day sequences. From our ACML 2026 study~\cite{ctde2026acml}, raw-OHLCVA LSTM achieves Calmar 6.05 vs. MLP on 53 features (Calmar 6.19) with 89\% fewer inputs. \textbf{BiLSTM} adds bidirectional context; \textbf{LSTM+Attention} uses learned temporal pooling.

\subsection{Bayesian-Optimized Tree Models}

\textbf{LightGBM}~\cite{ke2017lightgbm} and \textbf{XGBoost}~\cite{chen2016xgboost} operate on flattened OHLCVA windows (60$\times$6=360 features). Both are optimized via Optuna TPE with 9-parameter search.

\subsection{Multi-Agent Reinforcement Learning}

\textbf{CTDE-MARL}~\cite{ctde2026acml,lowe2017multi} uses FactorEncoder$\to$ConfidenceGate($\tau{=}0.5$)$\to$StrategyNet with CentralizedCritic. From our ablation~\cite{ctde2026acml}, the centralized critic accounts for $\sim$50\% of Calmar gain; the gate+IC-loss pair prevents strategy collapse.

\begin{table}[htbp]
\centering
\caption{Model zoo summary. $\dagger$: tested locally; $\ddagger$: Bayesian implementation complete, experiments pending.}
\label{tab:models}
\small
\begin{tabular}{@{}llrl@{}}
\toprule
\textbf{Family} & \textbf{Models} & \textbf{Params} & \textbf{Status} \\
\midrule
Foundation & Kronos (mini/small/base)$\ddagger$ & 4.1M--102.3M & Code complete \\
Temporal Transformer & Autoformer$\ddagger$, Informer$\ddagger$, TimesNet$\ddagger$ & 2--3M & Code complete \\
 & PatchTST$\ddagger$, Transformer$\dagger$ & 132K--2M & 1 tested \\
Recurrent & LSTM$\dagger$, GRU$\dagger$, BiLSTM$\dagger$ & 10--18K & 3 tested \\
 & LSTM+Attn$\dagger$, BiGRU, Deep LSTM & 8--150K & 1 tested \\
Tree (Bayesian) & LGBM$\ddagger$, XGBoost$\ddagger$ & N/A & Code complete \\
MARL & CTDE-MARL (LSTM/GRU/MLP) & 30--50K & Code complete \\
Knowledge & CausalKG Event Detector & N/A & Code complete \\
Ensemble & IC-weighted, Stacking & --- & Code complete \\
\bottomrule
\end{tabular}
\end{table}

%=============================================================================
\section{Bayesian Hyperparameter Optimization}
%=============================================================================

We apply Tree-structured Parzen Estimator (TPE)~\cite{bergstra2011algorithms} via Optuna~\cite{akiba2019optuna} to 10 model families. Each search runs 50 trials (20 epochs/trial, validated on Rank IC). Search spaces are in Table~\ref{tab:bayesian}. \textbf{Note:} search spaces are implemented but GPU experiments are pending — results will be added in revision.

\begin{table*}[htbp]
\centering
\caption{TPE Bayesian search spaces (50 trials each, 20 epochs/trial, objective = validation Rank IC). All spaces implemented; GPU experiments pending.}
\label{tab:bayesian}
\footnotesize
\begin{tabular}{@{}llp{6.5cm}l@{}}
\toprule
\textbf{Model} & \textbf{\#P} & \textbf{Search Space} & \textbf{Trials} \\
\midrule
BayesianKronos & 8 & variant$\in$\{mini,small\}, pred\_h$\in$[64,512], lr$\in$[1e-5,5e-3], bs$\in$[32,256], lookback$\in$[20,120] & 50 \\
BayesianAutoformer & 9 & d\_model$\in$[64,512], L$\in$[1,5], heads$\in$\{4,8\}, d\_ff$\in$[256,2048], lr$\in$[1e-5,5e-3] & 50 \\
BayesianInformer & 9 & Same as Autoformer & 50 \\
BayesianTimesNet & 9 & d\_model$\in$[64,512], L$\in$[1,4], d\_ff$\in$[16,64], top\_k$\in$[3,8], lr$\in$[1e-5,5e-3] & 50 \\
BayesianPatchTST & 10 & patch$\in$[8,32], stride$\in$[4,16], d\_model$\in$[64,256], L$\in$[1,4], heads$\in$\{4,8\} & 50 \\
BayesianTransformer & 8 & d\_model$\in$[64,512], L$\in$[1,6], heads$\in$\{4,8\}, d\_ff$\in$[256,2048], lr$\in$[1e-5,5e-3] & 50 \\
BayesianLSTM & 7 & hidden$\in$[32,256], L$\in$[1,4], drop$\in$[0,0.5], lr$\in$[1e-4,1e-2], bidir$\in$\{T,F\} & 50 \\
BayesianGRU & 7 & Same as LSTM & 50 \\
BayesianLGBM & 9 & n\_est$\in$[100,2000], depth$\in$[3,15], lr$\in$[0.01,0.3], leaves$\in$[15,255], $\alpha$/$\lambda$$\in$[1e-8,10] & 50 \\
BayesianXGBoost & 9 & n\_est$\in$[100,2000], depth$\in$[3,15], lr$\in$[0.01,0.3], $\gamma$$\in$[1e-8,5], $\alpha$/$\lambda$$\in$[1e-8,10] & 50 \\
\bottomrule
\end{tabular}
\end{table*}

%=============================================================================
\section{Local Baseline Experiment}
%=============================================================================

\subsection{Setup}
We validate the pipeline locally before BigQuant deployment. \textbf{Important caveat:} this is a \emph{smoke test}, not a competitive evaluation. The test uses:
\begin{itemize}
  \item \textbf{Data}: 10 A+H stocks, daily OHLCVA, 60-day lookback, Jan 2025--Apr 2026
  \item \textbf{Samples}: 2,146 train / 300 test
  \item \textbf{Hardware}: CPU (Intel i7)
  \item \textbf{Training}: 20 epochs, batch 128, lr $10^{-3}$, Adam, MSE loss, seed 42
  \item \textbf{Hyperparameters}: All defaults (no Bayesian optimization)
\end{itemize}

\subsection{Results}

Table~\ref{tab:local} reports results. All models exhibit near-zero cross-sectional IC. MLP achieves marginally positive IC (+0.038) due to lower variance under data scarcity.

\begin{table}[htbp]
\centering
\caption{Local baseline (CPU, daily, 10 stocks, 20 epochs, single seed=42).}
\label{tab:local}
\begin{tabular}{@{}lrrrrr@{}}
\toprule
\textbf{Model} & \textbf{Params} & \textbf{MSE} & \textbf{Rank IC} & \textbf{ICIR} & \textbf{Time (s)} \\
\midrule
MLP (baseline)   & 54,529 & 0.0020 & \textbf{+0.038} & +0.43 & 58.1 \\
BiLSTM           & 10,305 & 0.0011 & +0.008 & +0.60 & 48.9 \\
Transformer      & 132,481 & 0.0010 & --0.055 & +0.12 & 140.0 \\
GRU              & 13,889 & 0.0011 & --0.061 & --1.40 & 37.7 \\
LSTM+Attention   & 18,562 & 0.0010 & --0.093 & --0.37 & 32.0 \\
LSTM             & 18,497 & 0.0011 & --0.138 & --1.29 & 29.7 \\
\midrule
IC-Ensemble      & --- & 0.0010 & --0.022 & --- & --- \\
\bottomrule
\end{tabular}
\end{table}

\subsection{Analysis}

The near-zero IC is expected: daily OHLCVA with 2,146 samples provides minimal cross-sectional signal. MLP's marginal positive IC reflects low variance under data scarcity. Recurrent models exhibit negative IC from overfitting to noise. The Transformer (132K params) is severely underdetermined (16 samples/param). This confirms three necessary conditions for competition success: (1) \textbf{minute-frequency data} (480$\times$ observations), (2) \textbf{1000-stock cross-section} (100$\times$ for IC), and (3) \textbf{Bayesian HPO} to escape defaults.

\paragraph{Statistical note:} Single-seed results are insufficient for ranking. Multi-seed experiments (3+ seeds) will be reported with GPU results.

%=============================================================================
\section{Competition Strategy}
%=============================================================================

We propose a three-phase strategy (Table~\ref{tab:strategy}):

\begin{table}[htbp]
\centering
\caption{Three-phase submission strategy.}
\label{tab:strategy}
\begin{tabular}{@{}lllll@{}}
\toprule
\textbf{Phase} & \textbf{Dates} & \textbf{Models} & \textbf{Target IC} & \textbf{Sub/day} \\
\midrule
P1: Baseline    & Jul 1--7   & LSTM, GRU, BayesianLGBM/XGBoost & $>$0.03 & 3 \\
P2: Bayesian    & Jul 8--21  & Kronos, Autoformer, Informer, TimesNet, PatchTST & $>$0.05 & 3 \\
P3: Ensemble    & Jul 22--Aug 5 & IC-Weighted Top-5 Ensemble & $>$0.08 & 1--2 \\
\bottomrule
\end{tabular}
\end{table}

\textbf{Phase 1} establishes baselines with fast-to-train models at maximum submission frequency. \textbf{Phase 2} deploys 10 Bayesian-optimized deep models, each undergoing 50 TPE trials on GPU. \textbf{Phase 3} combines best models via IC-weighted ensemble ($w_i \propto e^{|\text{IC}_i|}$). Private leaderboard: best public model retrained from scratch in isolated environment.

%=============================================================================
\section{Limitations}
%=============================================================================

\begin{enumerate}
  \item \textbf{Experimental incompleteness:} Only 6 of 23 models tested; Bayesian optimization experiments are pending GPU access. Current results are smoke tests, not competitive evaluations.
  \item \textbf{Data constraints:} Daily-frequency data is insufficient for cross-sectional prediction. Minute-frequency data on BigQuant platform is required for meaningful IC.
  \item \textbf{Statistical rigor:} Single-seed results lack error bars. Multi-seed experiments ($\ge 3$ seeds) are planned.
  \item \textbf{Missing baselines:} iTransformer, DLinear, Mamba, and TimesFM are not yet included.
  \item \textbf{Task specificity:} Results are specific to BigAlpha 2026 competition constraints and may not generalize to other prediction tasks.
  \item \textbf{No transaction costs:} Current evaluation ignores market impact and execution costs.
\end{enumerate}

%=============================================================================
\section{Conclusion}
%=============================================================================

We have built a 23-model benchmark infrastructure for the BigAlpha 2026 competition, spanning foundation models, time-series transformers, recurrent architectures, Bayesian-optimized trees, and multi-agent RL. TPE-based Bayesian optimization covers 10 families with 7--10 parameter spaces. Local smoke testing confirms the necessity of GPU and minute-frequency data. A three-phase strategy targets IC $>$ 0.08 through systematic optimization and IC-weighted ensembling. Code is available at the project repository.

%=============================================================================
\bibliographystyle{ACM-Reference-Format}
\begin{thebibliography}{20}

\bibitem{kronos2025}
S.~Shi et al.
\newblock Kronos: A Foundation Model for the Language of Financial Markets.
\newblock \emph{AAAI}, 2026.

\bibitem{wu2021autoformer}
H.~Wu, J.~Xu, J.~Wang, and M.~Long.
\newblock Autoformer: Decomposition Transformers with Auto-Correlation.
\newblock \emph{NeurIPS}, 2021.

\bibitem{zhou2021informer}
H.~Zhou et al.
\newblock Informer: Beyond Efficient Transformer for Long Sequence Time-Series Forecasting.
\newblock \emph{AAAI}, 2021.

\bibitem{wu2023timesnet}
H.~Wu, T.~Hu, Y.~Liu, H.~Zhou, J.~Wang, and M.~Long.
\newblock TimesNet: Temporal 2D-Variation Modeling for General Time Series Analysis.
\newblock \emph{ICLR}, 2023.

\bibitem{nie2023patchtst}
Y.~Nie, N.~H. Nguyen, P.~Sinthong, and J.~Kalagnanam.
\newblock A Time Series is Worth 64 Words: Long-term Forecasting with Transformers.
\newblock \emph{ICLR}, 2023.

\bibitem{vaswani2017attention}
A.~Vaswani et al.
\newblock Attention Is All You Need.
\newblock \emph{NeurIPS}, 2017.

\bibitem{ke2017lightgbm}
G.~Ke et al.
\newblock LightGBM: A Highly Efficient Gradient Boosting Decision Tree.
\newblock \emph{NeurIPS}, 2017.

\bibitem{chen2016xgboost}
T.~Chen and C.~Guestrin.
\newblock XGBoost: A Scalable Tree Boosting System.
\newblock \emph{KDD}, 2016.

\bibitem{lowe2017multi}
R.~Lowe, Y.~Wu, A.~Tamar, J.~Harb, P.~Abbeel, and I.~Mordatch.
\newblock Multi-Agent Actor-Critic for Mixed Cooperative-Competitive Environments.
\newblock \emph{NeurIPS}, 2017.

\bibitem{ctde2026acml}
Z.~Li et al.
\newblock What Matters in Multi-Agent Reinforcement Learning for Portfolio Management?
\newblock \emph{ACML}, 2026. Under review.

\bibitem{bergstra2011algorithms}
J.~Bergstra, R.~Bardenet, Y.~Bengio, and B.~K\'{e}gl.
\newblock Algorithms for Hyper-Parameter Optimization.
\newblock \emph{NeurIPS}, 2011.

\bibitem{snoek2012practical}
J.~Snoek, H.~Larochelle, and R.~P. Adams.
\newblock Practical Bayesian Optimization of Machine Learning Algorithms.
\newblock \emph{NeurIPS}, 2012.

\bibitem{akiba2019optuna}
T.~Akiba, S.~Sano, T.~Yanase, T.~Ohta, and M.~Koyama.
\newblock Optuna: A Next-generation Hyperparameter Optimization Framework.
\newblock \emph{KDD}, 2019.

\bibitem{hochreiter1997long}
S.~Hochreiter and J.~Schmidhuber.
\newblock Long Short-Term Memory.
\newblock \emph{Neural Computation}, 9(8):1735--1780, 1997.

\bibitem{cho2014learning}
K.~Cho et al.
\newblock Learning Phrase Representations using RNN Encoder-Decoder.
\newblock \emph{EMNLP}, 2014.

\end{thebibliography}

\end{document}
